\nonstopmode
\documentclass[a4paper,11pt]{article}

\usepackage{fontspec}
\usepackage[a4paper, top=2.2cm, bottom=2.2cm, left=2.2cm, right=2.2cm]{geometry}
\usepackage{graphicx}
\usepackage{xcolor}
\usepackage{titlesec}
\usepackage{tabularx}
\usepackage{booktabs}
\usepackage{microtype}
\usepackage{hyperref}
\usepackage{enumitem}
\usepackage{array}
\usepackage{ragged2e}
\usepackage{parskip}
\usepackage{fancyhdr}
\usepackage{lastpage}
\usepackage{calc}
\usepackage{multirow}
\usepackage{mdframed}

\setmainfont[Ligatures=TeX]{Helvetica Neue}

\definecolor{raaRed}{HTML}{8B0000}
\definecolor{raaLight}{HTML}{f5f5dc}
\definecolor{raaMid}{HTML}{888888}
\definecolor{raaLine}{HTML}{cccccc}

\hypersetup{colorlinks=true, urlcolor=raaRed, linkcolor=raaRed}

\pagestyle{fancy}
\fancyhf{}
\renewcommand{\headrulewidth}{0pt}
\fancyfoot[L]{\footnotesize\color{raaMid} RAA.SPACE --- Matīsa iela 8, Rīga}
\fancyfoot[R]{\footnotesize\color{raaMid} \thepage\ / \pageref{LastPage}}

\titleformat{\section}{\large\bfseries\color{raaRed}}{}{0em}{}[\color{raaLine}\titlerule]
\titlespacing{\section}{0pt}{1.4em}{0.5em}

\titleformat{\subsection}{\normalsize\bfseries}{}{0em}{}
\titlespacing{\subsection}{0pt}{0.8em}{0.2em}

\newcommand{\evententry}[5]{%
  \noindent\begin{minipage}[t]{2.6cm}
    \vspace{0pt}
    \includegraphics[width=2.5cm,height=3cm,keepaspectratio]{#2}
  \end{minipage}%
  \hspace{0.3cm}%
  \begin{minipage}[t]{\linewidth - 2.9cm}
    \vspace{0pt}
    {\footnotesize\color{raaRed}\textbf{#1} \hfill \textit{\color{raaMid}#5}}\\[-0.1em]
    {\bfseries #3}\\[0.1em]
    {\small\RaggedRight #4}
  \end{minipage}%
  \par\vspace{0.4em}
  \noindent{\color{raaLine}\rule{\linewidth}{0.4pt}}
  \par\vspace{0.4em}
}

\newcommand{\milestone}[2]{%
  \noindent{\color{raaRed}\textbf{#1}} \quad #2\par\vspace{0.3em}
}

\newcommand{\teammember}[3]{%
  \begin{minipage}[t]{2.2cm}
    \vspace{0pt}
    \includegraphics[width=2cm,height=2.5cm,keepaspectratio]{#3}
  \end{minipage}%
  \hspace{0.3cm}%
  \begin{minipage}[t]{5.8cm}
    \vspace{0pt}
    {\bfseries #1}\\
    {\small\color{raaMid}#2}
  \end{minipage}%
}

\begin{document}

% ═════════════════════════════════════════════════════════
%  1. LAPPUSE — TITULLAPA + PAR RAA
% ═════════════════════════════════════════════════════════
\thispagestyle{empty}

\begin{center}
  \includegraphics[width=5cm,keepaspectratio]{/Users/leo/Projects/Raa/web/frontend/public/raa_logo_but_use_svg_if_possible.png}
  \par\vspace{0.6em}
  {\large\bfseries RAA.SPACE --- Kultūrtelpa}
  \par\vspace{0.2em}
  {\normalsize Radošā rūpnīca Veldze, Matīsa iela 8, Rīga LV-1001}
\end{center}

\vspace{0.4em}
\begin{center}
\includegraphics[width=\linewidth,height=8cm,keepaspectratio]{/Users/leo/Projects/Raa/web/frontend/public/assets/telpa/raa-galerija-1.jpg}
\end{center}
\vspace{0.4em}

\begin{center}
  \noindent{\color{raaLine}\rule{0.6\linewidth}{0.8pt}}
  \par\vspace{0.4em}
  \begin{tabular}{ccc}
    \textbf{E-pasts} & \textbf{Instagram} & \textbf{Adrese} \\[0.3em]
    \href{mailto:enter@raa.space}{enter@raa.space} &
    \href{https://instagram.com/raa_riga}{@raa\_riga} &
    Matīsa iela 8, Rīga LV-1001 \\
  \end{tabular}
  \par\vspace{0.4em}
  \noindent{\color{raaLine}\rule{0.6\linewidth}{0.8pt}}
  \par\vspace{0.3em}
  {\small\color{raaMid} Atsauces dokuments --- \today}
\end{center}

\vspace{0.5em}

% ─────────────────────────────────────────────────────────
\section{Par RAA}
% ─────────────────────────────────────────────────────────

RAA ir neatkarīga kultūrtelpa Radošajā rūpnīcā Veldze Rīgas centrā, Latvijā.
Tā atvērās \textbf{2025. gada 11. jūlijā}.

Telpa piedāvā daudzveidīgu programmu: laikmetīgās mākslas izstādes, teātra izrādes, koncertus,
dzejas lasījumus, darbnīcas, festivālus un kopienas pasākumus.

Telpa ietver izstāžu zāli (103\,m$^2$), ieejas vestibili (35\,m$^2$), teātra telpu
(13\,m$^2$) ar atsevišķu vestibili, aizkulišu koridoru un pieeju pagalmam, kas savieno
visu Veldzes kvartālu.

RAA (jeb \textit{reductio ad absurdum}) ved apgalvojumus līdz absurdam kopš 2025. gada.
RAA tev nav jābūt skaļam, lai tevi dzirdētu.

\vfill

\noindent{\color{raaLine}\rule{\linewidth}{0.4pt}}
{\footnotesize\color{raaMid}
  \textbf{RAA.SPACE} \enspace
  Biedrība \enspace
  Reģ.\ 40008349569 \enspace
  Matīsa iela 8, Rīga LV-1001
}

% ═════════════════════════════════════════════════════════
%  2. LAPPUSE — KOMANDA + MEŽA BĀRS + PLĀNOJUMS
% ═════════════════════════════════════════════════════════
\newpage

% ─────────────────────────────────────────────────────────
\section{Komanda}
% ─────────────────────────────────────────────────────────

\vspace{0.5em}
\noindent
\teammember{Leonhard Horstmeyer}{haosa meistars}{/Users/leo/Projects/Raa/web/frontend/public/assets/team/Leo.jpeg}%
\hspace{0.6cm}
\teammember{Jana Zariņa}{mākslinieciskā direktore}{/Users/leo/Projects/Raa/web/frontend/public/assets/team/JanaZarina.JPG}

\vspace{0.8em}
\noindent
\teammember{Helias Dzimoulakis}{teātra direktors}{/Users/leo/Projects/Raa/web/frontend/public/assets/team/HeliasD.jpg}%
\hspace{0.6cm}
\teammember{Toms Dāvis Piņķis}{skaņa, gaisma un bārs}{/Users/leo/Projects/Raa/web/frontend/public/assets/team/Toms-portrait.JPG}

\vspace{0.8em}
\noindent
\teammember{Madara Kurzemniece}{bārs un telpa}{/Users/leo/Projects/Raa/web/frontend/public/assets/team/Madara.jpeg}

% ─────────────────────────────────────────────────────────
\section{MAATE NAATRE}
% ─────────────────────────────────────────────────────────

MAATE NAATRE ir meža bārs RAA telpas ieejas vestibilā, Radošajā rūpnīcā Veldze Rīgas centrā.
Atklāts 2026. gada 14. martā, tas ir neliela vieta, kur var aizkavēties, satikties ar citiem
un nogaršot dažādus dabas radītus dzērienus. Bārs atnes Latvijas mežu uz pilsētu ---
izmēģini dažādus kokteiļus no meža.

% ─────────────────────────────────────────────────────────
\section{Plānojums}
% ─────────────────────────────────────────────────────────

\vspace{0.5em}
\noindent
\begin{minipage}[t]{0.495\linewidth}
  \vspace{0pt}
  \includegraphics[width=\linewidth,height=7cm,keepaspectratio]{/Users/leo/Projects/Raa/web/frontend/public/assets/telpa/raa-galerija-0.jpg}
  \par\vspace{0.3em}
  {\small\color{raaMid}\textit{Galvenā zāle}}
\end{minipage}%
\hfill%
\begin{minipage}[t]{0.495\linewidth}
  \vspace{0pt}
  \includegraphics[width=\linewidth,height=7cm,keepaspectratio]{/Users/leo/Projects/Raa/web/frontend/public/assets/telpa/layout.jpeg}
  \par\vspace{0.3em}
  {\small\color{raaMid}\textit{RAA plānojums (bez sarkanās zonas)}}
\end{minipage}

% ═════════════════════════════════════════════════════════
%  3. LAPPUSE — HRONOLOĢIJA
% ═════════════════════════════════════════════════════════
\newpage

% ─────────────────────────────────────────────────────────
\section{Hronoloģija}
% ─────────────────────────────────────────────────────────

\milestone{Jūlijs 2025}{RAA atvērās 2025. gada 11. jūlijā. Atklāšanas grupu izstāde
  \textit{In Your Shoes}, kuratorēja Maija Kurševa, pulcēja septiņus māksliniekus no
  Latvijas Mākslas akadēmijas. Dzejas maratons \textit{Sinusoīda} ar Syntax Stupor.}

\milestone{Augusts 2025}{Katoptron Teātris, RAA rezidenta teātris, atvērās Helia Dzimoulaka
  vadībā. Arta Raituma vadīja divu dienu uguns darbnīcu, kam sekoja radīto darbu izstāde.
  Frīča Ripas soloizstāde \textit{Saules Gaismas Meklējumos}. Muzikāls vakars ar Elzu
  un Uģi Albiņiem. Ksenijas Kumačovas fotoizstāde.}

\milestone{Septembris 2025}{Homo Novus festivālam veltīta izstāde. Šujmašīnu tehno.
  Lielbritānijas dzejnieces Džoelas Teilores (T. S. Eliota balva) dzejas lasījums.
  Tvēruma \#11 numura prezentācija. Vakars ar Eriku Kolobanova un Konstantīnu Belihu.}

\milestone{Oktobris 2025}{Trīs dienu Starptautiskais smaržu mākslas festivāls (10.--12.\ okt.)
  ar māksliniekiem no dažādām valstīm. res Tacta performances un darbnīcu sērija.
  Brīvības Experience (16.--30.\ okt.) --- pilsētvides instalācija par Brīvības ielu,
  ko organizē Brunis Kibilds (Pilsēta cilvēkiem).}

\milestone{Novembris 2025}{Kefirs Studio grupu izstāde. Pirmais Rīgas Upcycling festivāls
  (15.\ nov.). Veronika Geidas fotoizstāde. \textit{Through the Looking Glass} ---
  Paulas Betijas Antenišķes fotogrāfija, dokumentējot Ivo Krieviņu un Rīgas ballroom scēnu.}

\milestone{Decembris 2025}{WunderKombināts publikācijas prezentācija. Klavierkoncerts ar
  Andreju Osokinsu. Tvēruma decembra numura prezentācija.}

\milestone{Janvāris 2026}{Ziemas pauze un renovācijas darbi RAA.}

\milestone{Februāris 2026}{30 dienu attīstoša instalācija \textit{AS DISTANT, SO NEAR} ar
  dažādiem māksliniekiem, kas pēta attāluma, tuvuma un klātbūtnes tēmas.}

\milestone{Marts 2026}{MAATE NAATRE meža bāra atklāšana (14.\ marts). Trīs pētniecības
  sanāksmju sesijas. Grāmatas prezentācija. Audēju grupu izstāde un Leonharda Horstmeyera
  performance \textit{Dialogo ad Infinitum}.}

\milestone{Aprīlis 2026}{\textit{ALTERNATĪVA: Pēdas no pazemes} --- 2ANNAS festivāla
  izstāde (10.\ apr. -- 4.\ maijs), ar četriem māksliniekiem un trim ekskursijām,
  ko atbalsta Rīgas dome un VKKF.}

% ═════════════════════════════════════════════════════════
%  4. LAPPUSE+ — PASĀKUMI
% ═════════════════════════════════════════════════════════
\newpage

% ─────────────────────────────────────────────────────────
\section{Pasākumi}
% ─────────────────────────────────────────────────────────

{\small Visi RAA rīkotie pasākumi no atklāšanas 2025. gada 11. jūlijā hronoloģiskā secībā.
2026. gadam plānoti arī daudzi citi pasākumi, kas vēl nav iekļauti šajā sarakstā.}

\vspace{0.6em}

% ── JŪLIJS 2025 ──────────────────────────────────────────
\subsection*{\color{raaMid} Jūlijs 2025}

\evententry%
  {11. jūlijs 2025}%
  {/Users/leo/Projects/Raa/web/frontend/public/assets/raa-opening/banner-opening-thumb.jpeg}%
  {R A A Opening \#1}%
  {RAA kultūrtelpas atklāšana Matīsa ielā 8.}%
  {Atklāšana}

\evententry%
  {11.--18. jūlijs 2025}%
  {/Users/leo/Projects/Raa/web/frontend/public/assets/tavaas-kurpees/poster-thumb.jpg}%
  {In Your Shoes}%
  {Grupu izstāde, kuratorē Maija Kurševa (LMA KAS 3. kursa studenti).
  Katrs darbs atsaucas uz cita mākslinieka praksi. Septiņi mākslinieki.}%
  {Izstāde}

\evententry%
  {24.--27. jūlijs 2025}%
  {/Users/leo/Projects/Raa/web/frontend/public/assets/typewriter-museum/poster-thumb.jpg}%
  {Typewriter Museum Exhibition}%
  {20. gs. rakstāmmašīnu kolekcija no Vācijas, ASV, Zviedrijas un Itālijas.
  Apmeklētāji varēja rakstīt uz divām autentiskām mašīnām.}%
  {Izstāde}

\evententry%
  {27. jūlijs 2025}%
  {/Users/leo/Projects/Raa/web/frontend/public/assets/dzeja-sinusoiida/RAA_dzejasSinusidaPRINT-thumb.jpg}%
  {Dzeja: Sinusoīda}%
  {Dzejas maratons ar grupu Syntax Stupor, mūziku spēlē Dabsolis un Kaspars Sils.
  Atvērts mikrofons.}%
  {Dzeja / Mūzika}

% ── AUGUSTS 2025 ────────────────────────────────────────
\subsection*{\color{raaMid} Augusts 2025}

\evententry%
  {2.--3. augusts 2025}%
  {/Users/leo/Projects/Raa/web/frontend/public/assets/fire-workshop/uguns_poster_DRUKA_with_QR_code-thumb.png}%
  {Uguns darbnīca ar Artu Raitumu}%
  {Divu dienu eksperimentāla koka dedzināšanas darbnīca ar mākslinieci Artu Raitumu.}%
  {Darbnīca}

\evententry%
  {4.--6. augusts 2025}%
  {/Users/leo/Projects/Raa/web/frontend/public/assets/fire-workshop/uguns_poster_DRUKA_with_QR_code-thumb.png}%
  {Uguns darbnīcas izstāde}%
  {Darbnīcā radīto darbu izstāde.}%
  {Izstāde}

\evententry%
  {7.--10. augusts 2025}%
  {/Users/leo/Projects/Raa/web/frontend/public/assets/kas-aug-starp-mums/poster-thumb.png}%
  {Kas Aug Starp Mums}%
  {Ksenijas Kumačovas (dz. 1997) fotoizstāde.
  Pēta cilvēka un vietējās dabas mijiedarbību.}%
  {Izstāde}

\evententry%
  {9. augusts 2025}%
  {/Users/leo/Projects/Raa/web/frontend/public/assets/theatre-opening/poster-thumb.png}%
  {Katoptron Teātris atklāšana}%
  {RAA rezidenta teātra Katoptron Teātris atklāšana. Dibinājis un vada Helias Dzimoulakis.}%
  {Teātris}

\evententry%
  {14. augusts 2025}%
  {/Users/leo/Projects/Raa/web/frontend/public/assets/kabiirijas-salons/poster-thumb.jpg}%
  {Kabiirijas Salons \#1}%
  {Atvērta skatuve glezniecībai, mūzikai, performancei, dejai un dzejai.}%
  {Performance / Atvērtā skatuve}

\evententry%
  {14.--17. augusts 2025}%
  {/Users/leo/Projects/Raa/web/frontend/public/assets/shelf-stories/impression-thumb.jpeg}%
  {RAA Muzejs: Plauktu Stāsti}%
  {Muzejs ar priekšmetiem, ko kopienas locekļi ziedojuši ar saviem stāstiem.}%
  {Izstāde}

\evententry%
  {18.--26. augusts 2025}%
  {/Users/leo/Projects/Raa/web/frontend/public/assets/saules-gaismas-mekleejumos/poster1-thumb.jpeg}%
  {Saules Gaismas Meklējumos}%
  {Frīča Ripas soloizstāde. Atklāšanā DJ seti.}%
  {Izstāde}

\evententry%
  {28. augusts 2025}%
  {/Users/leo/Projects/Raa/web/frontend/public/assets/laime-pliist-uz-traukiem/poster-thumb.jpg}%
  {Laime Plīst uz Traukiem}%
  {Vakara uzveduma pasākums.}%
  {Performance}

% ── SEPTEMBRIS 2025 ─────────────────────────────────────
\subsection*{\color{raaMid} Septembris 2025}

\evententry%
  {1. septembris 2025}%
  {/Users/leo/Projects/Raa/web/frontend/public/assets/second-variety-show/poster-thumb.jpg}%
  {Second Variety Show}%
  {Dažādu priekšnesumu vakars.}%
  {Performance}

\evententry%
  {3. septembris 2025}%
  {/Users/leo/Projects/Raa/web/frontend/public/assets/pilseetas-naktii/poster-thumb.jpg}%
  {Pilsētas Naktī}%
  {Vakara pasākums par pilsētvides tematiku.}%
  {Pasākums}

\evententry%
  {8. septembris 2025}%
  {/Users/leo/Projects/Raa/web/frontend/public/assets/dzoela-teilore/Dzoela_Teilore_Raa_A3-thumb.jpg}%
  {Džoela Teilore --- dzejas lasījums}%
  {Dzejas lasījums ar Lielbritānijas dzejnieci Džoelu Teilori, T. S. Eliota balvas laureāti.}%
  {Literatūra / Dzeja}

\evententry%
  {14. septembris 2025}%
  {/Users/leo/Projects/Raa/web/frontend/public/assets/sujmasinu-tehno/poster-thumb.png}%
  {Šujmašīnu Tehno}%
  {Vakars, kas apvieno amatniecību, kopienu un ritmu.}%
  {Darbnīca / Pasākums}

\evententry%
  {18. septembris 2025}%
  {/Users/leo/Projects/Raa/web/frontend/public/assets/ja-nu-saule-rit-neuzaust/poster-thumb.jpeg}%
  {Tvērums \#11 ``Austrumi'': Ko darīt, ja saule rīt neuzaust}%
  {Literārā žurnāla Tvērums 11. numura prezentācija.}%
  {Literatūra / Prezentācija}

\evententry%
  {27. septembris 2025}%
  {/Users/leo/Projects/Raa/web/frontend/public/assets/exo-foto-izstade/poster-thumb.jpeg}%
  {Exo --- fotoizstāde}%
  {Fotoizstāde.}%
  {Izstāde}

% ── OKTOBRIS 2025 ───────────────────────────────────────
\subsection*{\color{raaMid} Oktobris 2025}

\evententry%
  {4. oktobris 2025}%
  {/Users/leo/Projects/Raa/web/frontend/public/assets/kizomba/poster-thumb.png}%
  {Kizomba / UrbanKiz ballīte}%
  {Sociālās dejas ballīte ar Kizomba un UrbanKiz mūziku.}%
  {Deja / Ballīte}

\evententry%
  {5. oktobris 2025}%
  {/Users/leo/Projects/Raa/web/frontend/public/assets/riga-sock-exchange/poster-thumb.png}%
  {Rīgas zeķu maiņa}%
  {Kopienas apģērbu maiņas pasākums.}%
  {Kopiena}

\evententry%
  {8. oktobris 2025}%
  {/Users/leo/Projects/Raa/web/frontend/public/assets/res-tacta/poster-thumb.png}%
  {res tacta --- performance / darbnīca}%
  {Performance un darbnīca, kas pēta taustes pieredzi.}%
  {Performance / Darbnīca}

\evententry%
  {10.--12. oktobris 2025}%
  {/Users/leo/Projects/Raa/web/frontend/public/assets/scent-art-festival/poster-thumb.jpg}%
  {Starptautiskais smaržu mākslas festivāls}%
  {Trīs dienu starptautisks festivāls, kas veltīts ožas mākslinieciskajai izpausmei ---
  izstādes, performances, darbnīcas un smaržu ekskursija pa Rīgu.}%
  {Festivāls}

\evententry%
  {16.--30. oktobris 2025}%
  {/Users/leo/Projects/Raa/web/frontend/public/assets/brivibas-experience/poster-thumb.png}%
  {Brīvības Experience --- pilsētvides instalācija}%
  {Izstāde un pilsētvides instalācija, kas veltīta Rīgas galvenajai ielai, Brīvības ielai.
  Izmantojot vizualizācijas, 3D modeļus un citus medijus, projekts pēta ielas vēsturi,
  pašreizējo stāvokli un nākotnes scenārijus. Organizators: Brunis Kibilds
  (Pilsēta cilvēkiem). Apmeklētāji tika aicināti piedalīties diskusijā par
  cilvēkdraudzīgāku pilsētvidi.}%
  {Pasākums}

% ── NOVEMBRIS 2025 ──────────────────────────────────────
\subsection*{\color{raaMid} Novembris 2025}

\evententry%
  {6.--9. novembris 2025}%
  {/Users/leo/Projects/Raa/web/frontend/public/assets/kefirs-studio-exhibition/poster-thumb.jpg}%
  {Kefirs Studio izstāde}%
  {Kefirs Studio grupu izstāde.}%
  {Izstāde}

\evententry%
  {15. novembris 2025}%
  {/Users/leo/Projects/Raa/web/frontend/public/assets/pirmais-rigas-upcycling-festivals/poster-thumb.jpg}%
  {Pirmais Rīgas Upcycling festivāls}%
  {Publisks pasākums, kas veicina ilgtspējīgu dizainu un otrreizēju pārstrādi.
  Darbnīcas, tirgus un kopienas aktivitātes.}%
  {Festivāls}

\evententry%
  {17. novembris 2025}%
  {/Users/leo/Projects/Raa/web/frontend/public/assets/pa-liekas-izstade/poster-thumb.jpeg}%
  {(pa)liekas --- fotoizstāde}%
  {Veronika Geidas soloizstāde.}%
  {Izstāde}

\evententry%
  {24.--30. novembris 2025}%
  {/Users/leo/Projects/Raa/web/frontend/public/assets/through-the-looking-glass/poster-thumb.jpg}%
  {Through the Looking Glass}%
  {Fotoizstāde, kas pēta kvīru identitāti un ballroom kultūru Latvijā.
  Fotogrāfija --- Paula Betija Antenišķe, dokumentējot Ivo Krieviņu un Rīgas ballroom scēnu.
  Atbalstīts ar Eiropas Savienības līdzfinansējumu.}%
  {Izstāde}

% ── DECEMBRIS 2025 ──────────────────────────────────────
\subsection*{\color{raaMid} Decembris 2025}

\evententry%
  {9. decembris 2025}%
  {/Users/leo/Projects/Raa/web/frontend/public/assets/wunderkombinats/poster-thumb.jpg}%
  {WunderKombināts --- prezentācija}%
  {WunderKombināts publikācijas prezentācija.}%
  {Literatūra / Prezentācija}

\evententry%
  {12. decembris 2025}%
  {/Users/leo/Projects/Raa/web/frontend/public/assets/andrejs-osokins-koncerts/poster-thumb.jpg}%
  {Andrejs Osokins --- klavierkoncerts}%
  {Klavierkoncerts ar Andreju Osokinsu.}%
  {Koncerts}

\evententry%
  {19. decembris 2025}%
  {/Users/leo/Projects/Raa/web/frontend/public/assets/tverums-decembris/poster-thumb.jpg}%
  {Tvērums --- decembra numura prezentācija}%
  {Tvēruma decembra numura prezentācija.}%
  {Literatūra / Prezentācija}

\evententry%
  {28. decembris 2025}%
  {/Users/leo/Projects/Raa/web/frontend/public/assets/marks-ozolins-koncerts/poster-thumb.jpg}%
  {Mārks Ozoliņš --- koncetvakars}%
  {Koncetvakars ar Mārku Ozoliņu.}%
  {Koncerts}

% ── JANVĀRIS 2026 ───────────────────────────────────────
\subsection*{\color{raaMid} Janvāris 2026}

\evententry%
  {29. janvāris 2026}%
  {/Users/leo/Projects/Raa/web/frontend/public/assets/as-far-so-near/poster-thumb.jpeg}%
  {AS DISTANT, SO NEAR --- atklātais konkurss}%
  {Atklātais konkurss un izstāde, kas pēta attāluma, tuvuma un klātbūtnes tēmas.
  30 dienu attīstošas instalācijas sākums ar dažādiem māksliniekiem.}%
  {Izstāde / Atklātais konkurss}

% ── MARTS 2026 ─────────────────────────────────────────
\subsection*{\color{raaMid} Marts 2026}

\evententry%
  {14. marts 2026}%
  {/Users/leo/Projects/Raa/web/frontend/public/assets/upcycling-tirgus-march/poster-thumb.jpg}%
  {2. Upcycling festivāls ``Atjauno veco!''}%
  {Otrais Rīgas Upcycling festivāls. Tirgus, darbnīcas un kopienas aktivitātes.}%
  {Festivāls}

\evententry%
  {14. marts 2026}%
  {/Users/leo/Projects/Raa/web/frontend/public/assets/maate-naatre-atklaashana/poster-thumb.jpg}%
  {MAATE NAATRE --- atklāšana}%
  {MAATE NAATRE meža bāra atklāšana RAA telpā.}%
  {Atklāšana}

\evententry%
  {17., 24.\ un 27. marts 2026}%
  {/Users/leo/Projects/Raa/web/frontend/public/assets/citizen-research-meetup-first-meeting/poster-thumb.png}%
  {Pētniecības sanāksmju sērija}%
  {Trīs atvērtas sanāksmes RAA zināšanu apmaiņai un kolektīvai izpētei.
  \textbf{Pilsoņu pētniecības sanāksme} (17.\ marts) --- atvērtu sesiju sākums datu
  jautājumu apspriešanai un jaunu prasmju apguvei. Organizē jopalesi.science.
  \textbf{Rekreatīvās pētniecības sanāksme} (24.\ marts) --- neformāla ideju
  un zinātkāru jautājumu apmaiņa.
  \textbf{Ārkārtas sanāksme: Rīmana dzeta} (27.\ marts) --- īpaša sesija par
  Rīmana dzeta funkciju un matemātisko skaistumu.}%
  {Sanāksme / Pētniecība}

\evententry%
  {26. marts 2026}%
  {/Users/leo/Projects/Raa/web/frontend/public/assets/weavers/poster-thumb.png}%
  {Audēji --- izstāde}%
  {Tekstila un starpnozaru izstāde.}%
  {Izstāde}

\evententry%
  {26. marts 2026}%
  {/Users/leo/Projects/Raa/web/frontend/public/assets/weavers/poster-thumb.png}%
  {Dialogo ad Infinitum --- performance}%
  {Leonharda Horstmeyera performance. Daļa no Audēju izstādes.}%
  {Performance}

\evententry%
  {28. marts 2026}%
  {/Users/leo/Projects/Raa/web/frontend/public/assets/laimes-govs/poster-thumb.jpg}%
  {Laimes Govs}%
  {Kopienas rīta pasākums.}%
  {Pasākums}

% ── APRĪLIS -- MAIJS 2026 ──────────────────────────────
\subsection*{\color{raaMid} Aprīlis -- Maijs 2026}

\evententry%
  {10. aprīlis -- 4. maijs 2026}%
  {/Users/leo/Projects/Raa/web/frontend/public/assets/alternativa-2-annas/poster-thumb.jpeg}%
  {ALTERNATĪVA: Pēdas no pazemes}%
  {2ANNAS festivāla izstāde. Mākslinieki: Laimdota Malle (LV), alliance Kurss (LV),
  Rebecca Korn (UK), Anastasija Šnepe-Šnepe (LV). Kurators: Kristiāna Kārkliņa.
  Trīs ekskursijas. Atbalsta Rīgas dome un VKKF.}%
  {Izstāde}

\evententry%
  {8.--10. maijs 2026}%
  {/Users/leo/Projects/Raa/web/frontend/public/assets/es-esmu-seit/poster-thumb.jpg}%
  {Atrodi mani, es esmu šeit}%
  {LMA studentu grupu izstāde kursa ``Izstāžu veidošanas prakse māksliniekiem'' ietvaros.
  Kuratorē Ieva Stalšene.}%
  {Izstāde}

\evententry%
  {15. maijs -- 8. jūnijs 2026}%
  {/Users/leo/Projects/Raa/web/frontend/public/assets/TP054K-01xG/poster-thumb.jpeg}%
  {TP054K-01xG}%
  {Dueta izstāde --- Alīna Burlakova (dz. 1992) un Viktorija Galkina (dz. 1988).
  Atklāšana 15. maijā ar Eleonoras Kampes vokālo performance.}%
  {Izstāde}

\end{document}
